%Master File:lectures.tex

\lesson{Neural Networks for Biological Prediction}
\vspace{-2cm}
\begin{center}
 \includegraphics[height=8cm]{Secondary_Protein_Structures}
 \includegraphics[height=8cm]{secondaryStructure}
\end{center}

\keywords{MLP, One-Hot Vectors, Softmax, Profile, Psi-Blast}
%%%%%%%%%%%%%%%%%%%%%%% Next Slide %%%%%%%%%%%%%%%%%%%%%%%
\renewcommand{\Outline}{%
\begin{slide}
\section[1]{Outline}

\begin{minipage}{14cm}
  \begin{enumerate}\squeeze
    \outlineitem{MLP}{mlp}
    \outlineitem{Secondary Protein Structure}{sps}
    \outlineitem{Designing MLPs}{dmlp}
   \end{enumerate}
\end{minipage}\hfill
\begin{minipage}{7cm}
  \vfill
  \includegraphics[height=12cm]{secondary-structure-of-protein}
  \vfill
\end{minipage}
\end{slide}
\addtocounter{outlineitem}{1}
}

\setcounter{outlineitem}{1}
%%%%%%%%%%%%%%%%%%%%%%% Next Slide %%%%%%%%%%%%%%%%%%%%%%%
\Outline
\toptarget{firstoutline}
%%%%%%%%%%%%%%%%%%%%%%% Next Slide %%%%%%%%%%%%%%%%%%%%%%%

\begin{slide}
\section{Neural Networks in Biology}

\begin{PauseHighLight}
  \begin{itemize}
  \item Neural networks (and support vector machines) are frequently
    used to make prediction in biology\pause
  \item We are not going into detail about neural networks (many of you
    have done this in machine learning)\pause
  \item However, how to use them in biology is quite subtle and is the
    main topic of this lecture\pause
  \item We take as a case study the classic problem of the prediction
    of secondary structure of proteins\pause
  \item Although, some of this is dated, the lessons are timeless\pause
  \end{itemize}
\end{PauseHighLight}

\end{slide}

%%%%%%%%%%%%%%%%%%%%%%% Next Slide %%%%%%%%%%%%%%%%%%%%%%%

\begin{slide}
\section[-2]{Multilayer Perceptron (MLP)}

\pause
\pb
\pauselevel{=1}
\begin{center}
  \multipdf[width=0.9\linewidth]{mlp}\pause
\end{center}
\end{slide}

%%%%%%%%%%%%%%%%%%%%%%% Next Slide %%%%%%%%%%%%%%%%%%%%%%%

\begin{slide}
\section{Training Neural Networks}

\begin{PauseHighLight}
  \begin{itemize}
  \item We train a neural network on a set of data
    \begin{align*}
      \mathcal{D} = \left( (\bm{x}^\mu, \bm{y}^\mu) \, \big\vert \, \mu
      = 1,\,2,\, \ldots,\, m\right)\pause
    \end{align*}
  \item Define an loss function $L(\bm{\theta}, \mathcal{D})$\pause
  \item E.g. squared error
    \begin{align*}
      L(\bm{\theta}, \mathcal{D}) = \frac{1}{m}
      \sum_{\mu=1}^m \left\| \bm{y}^\mu -
      \bm{f}(\bm{x}^\mu| \bm{\theta}) \right\|^2\pause
    \end{align*}
  \end{itemize}
\end{PauseHighLight}

\end{slide}

%%%%%%%%%%%%%%%%%%%%%%% Next Slide %%%%%%%%%%%%%%%%%%%%%%%

\begin{slide}
\section[-1]{Gradient Descent}

\begin{PauseHighLight}
  \begin{itemize}
  \item To minimise the loss we follow the negative gradient
    \begin{align*}
      \grad L(\bm{\theta}, \mathcal{D}) =
      \begin{pmatrix}
        \pd{L(\bm{\theta}, \mathcal{D})}{\theta_1} \\
        \pd{L(\bm{\theta}, \mathcal{D})}{\theta_2} \\
        \vdots \\
        \pd{L(\bm{\theta}, \mathcal{D})}{\theta_N}
      \end{pmatrix}\pause
    \end{align*}
  \item where $\theta_1=W_{11}$, $\theta_2=W_{12}$, etc.\pause
  \item We make a step in the direction of the negative gradient
    \begin{align*}
      \bm{\theta}(t+1) = \bm{\theta}(t) - r\,\grad L(\bm{\theta}, \mathcal{D})
    \end{align*}
  \item $r$ is learning rate\pause
  \end{itemize}
\end{PauseHighLight}

\end{slide}

%%%%%%%%%%%%%%%%%%%%%%% Next Slide %%%%%%%%%%%%%%%%%%%%%%%

\begin{slide}
\section{Following the Gradient}
\pause
\pb
\pauselevel{=1}
\begin{center}
  \multipdf[width=\linewidth]{gradientDescent}\pause
\end{center}

\end{slide}

%%%%%%%%%%%%%%%%%%%%%%% Next Slide %%%%%%%%%%%%%%%%%%%%%%%

\begin{slide}
\section{Backpropagation}

\begin{PauseHighLight}
  \begin{itemize}
  \item Calculating the gradient for a MLP is known as
    back-propagation\pause
  \item We can compute the gradient of the error for the whole training
    set\pause
  \item In this case we often use Levenberg-Marquardt, scale conjugate
    gradient or BFGS to speed up convergence\pause
  \item More often it is effective to find the error per input pattern
    or mini-batch (a small number of input patterns)\pause
  \item This goes by the name of stochastic gradient descent\pause
  \item This introduces noise which helps escape local optima\pause
  \end{itemize}
\end{PauseHighLight}

\end{slide}

%%%%%%%%%%%%%%%%%%%%%%% Next Slide %%%%%%%%%%%%%%%%%%%%%%%

\begin{slide}
\section{Higher Dimensions}

\pb
\begin{itemize}
\item In higher dimensions the problem is that there are some directions
  you need to move a long way\pauseh
\end{itemize}
\pauselevel{=1}
\begin{center}
  \includegraphics[width=\linewidth]{min2d-bg}\pause
  \multido{\ia=0+1}{10}{%
    \llap{\includegraphics[width=\linewidth]{min2d20-\ia}}\pause}
\end{center}
\end{slide}


%%%%%%%%%%%%%%%%%%%%%%% Next Slide %%%%%%%%%%%%%%%%%%%%%%%

\begin{slide}
\section{More Haste Less Speed}

\pb
\begin{itemize}
\item Increasing the step size a little  increases the
  rate of converge in one direction, but \ldots\pauseh
\end{itemize}
\begin{center}
  \includegraphics[width=\linewidth]{min2d-bg}\pause
  \multido{\ia=0+1}{10}{%
    \llap{\includegraphics[width=\linewidth]{min2d35-\ia}}\pause}
\end{center}

\end{slide}


%%%%%%%%%%%%%%%%%%%%%%% Next Slide %%%%%%%%%%%%%%%%%%%%%%%

\begin{slide}
\section{Speeding Up Stochastic Gradient Descent}

\begin{PauseHighLight}
  \begin{itemize}
  \item To speed up stochastic gradient descent we can add a momentum
    term
    \begin{align*}
      \bm{v}(t) &= (1-\epsilon)\bm{v}(t) - r \grad
                  L(\bm{\theta}(t-1))
                  \\
      \bm{\theta}(t) &= \bm{\theta}(t-1) + \bm{v}(t)\pause
    \end{align*}
  \item Many variations on this, e.g. ada-delta and adam\pause
  \item These are quite new techniques developed for deep-learning, but
    equally applicable to MLPs\pause
  \end{itemize}
\end{PauseHighLight}

\end{slide}

%%%%%%%%%%%%%%%%%%%%%%% Next Slide %%%%%%%%%%%%%%%%%%%%%%%
\Outline % Secondary Protein Structture
%%%%%%%%%%%%%%%%%%%%%%% Next Slide %%%%%%%%%%%%%%%%%%%%%%%

\begin{slide}
\section{Protein Structure}

\begin{PauseHighLight}
  \begin{itemize}
  \item Understanding the 3D structure of proteins is very important to
    understand its function\pause
  \item Unfortunately, it is extremely difficult to infer its 3D
    structure\pause
    \begin{itemize}
    \item It is both experimentally very costly\pause
    \item And computationally very challenging\pause
    \end{itemize}
  \item However knowing a proteins \emph{primary} structure (sequence)
    is quite easy (we can read this off if we know the RNA sequence from
    which it is translated)\pause
  \end{itemize}
\end{PauseHighLight}

\end{slide}

%%%%%%%%%%%%%%%%%%%%%%% Next Slide %%%%%%%%%%%%%%%%%%%%%%%

\begin{slide}
\section{Secondary Structure}

\begin{PauseHighLight}
  \begin{itemize}
  \item The secondary structure tells us whether a protein is part of
    \begin{itemize}
    \item An alpha-helix
    \item A beta-sheet
    \item Something else (often referred to as coil)\pause
    \end{itemize}
  \item This does not tell us the full 3-D structure but it is the first
    step towards inferring this\pause
  \end{itemize}
\end{PauseHighLight}

\end{slide}

%%%%%%%%%%%%%%%%%%%%%%% Next Slide %%%%%%%%%%%%%%%%%%%%%%%

\begin{slide}
\section[-2]{Protein Structure}

\begin{center}
  \includegraphics[height=15cm]{secondaryStructure}
\end{center}
\end{slide}

%%%%%%%%%%%%%%%%%%%%%%% Next Slide %%%%%%%%%%%%%%%%%%%%%%%

\begin{slide}
\section{Test Problem}

\begin{minipage}{0.5\linewidth}
  \begin{PauseHighLight}
  \begin{itemize}
  \item Secondary protein structure has long been a classic test problem
    for machine learning\pause
  \item It was an early success for MLPs\pause
  \item People have tried SVMs, HMMs, etc.\pause
  \item It is a difficult problem in part because it is non-local\pause
  \end{itemize}
  \end{PauseHighLight}
\end{minipage}\hfil
\begin{minipage}{0.48\linewidth}
  \begin{center}
    \includegraphics[width=\linewidth]{alphabeta}
  \end{center}
\end{minipage}

\end{slide}


%%%%%%%%%%%%%%%%%%%%%%% Next Slide %%%%%%%%%%%%%%%%%%%%%%%
\Outline % Desigining MLPs
%%%%%%%%%%%%%%%%%%%%%%% Next Slide %%%%%%%%%%%%%%%%%%%%%%%

\begin{slide}
\section{Predicting Secondary Structure}

\begin{PauseHighLight}
  \begin{itemize}
  \item Although very costly we know the secondary structure for quite a
    few proteins\pause
  \item This provides a training and testing data set\pause
  \item However, its not immediately obvious how to create a neural
    network to use this data\pause
  \item Most learning machines take a fixed length input\pause
  \end{itemize}
\end{PauseHighLight}

\end{slide}

%%%%%%%%%%%%%%%%%%%%%%% Next Slide %%%%%%%%%%%%%%%%%%%%%%%

\begin{slide}
\section{Windowing}
\pause
\pb
\pauselevel{=1}
\begin{center}
  \multipdf[width=\linewidth]{windowing}\pause
\end{center}

\end{slide}

%%%%%%%%%%%%%%%%%%%%%%% Next Slide %%%%%%%%%%%%%%%%%%%%%%%

\begin{slide}
\section{Randomised Windows}
\pause
\pb
\begin{itemize}
\item During training we can randomise the placement of the window so we
  don't learn a long section of one structure\pauseh
\pauselevel{=1}

\begin{center}
  \multipdf[width=\linewidth]{randomWindowing}\pause
\end{center}

\end{itemize}

\end{slide}


%%%%%%%%%%%%%%%%%%%%%%% Next Slide %%%%%%%%%%%%%%%%%%%%%%%

\begin{slide}
\section[-2]{Representing Proteins}

\begin{PauseHighLight}
  \begin{itemize}
  \item How should we feed our protein to an MLP?\pause
  \item Providing a number from 1--20 would put in an artificial
    ordering\pause
  \item We could try and put in chemical properties, e.g. acid/base
    hydophobic/hydophilic, etc.\pause{} but which ones\pause
  \item It turned out best to use a \emph{1-hot} (orthogonal) representation
    \begin{align*}
      \text{Glycine (G)} &= (1,0,0,0,0,0,0,0,0,0,0,0,0,0,0,0,0,0,0,0,0) \\
      \text{Alanine (A)} &= (0,1,0,0,0,0,0,0,0,0,0,0,0,0,0,0,0,0,0,0,0) \\
      \vdots\\
      \text{End} &= (0,0,0,0,0,0,0,0,0,0,0,0,0,0,0,0,0,0,0,0,1)\pause
    \end{align*}
  \item This is unbiased and seems to work best\pause
  \end{itemize}
\end{PauseHighLight}

\end{slide}

%%%%%%%%%%%%%%%%%%%%%%% Next Slide %%%%%%%%%%%%%%%%%%%%%%%

\begin{slide}
\section{Output}

\begin{PauseHighLight}
  \begin{itemize}
  \item Usually have 3 output (alpha-helix, beta-sheet and coil)
    although sometimes coil is split into 6 classes\pause
  \item Again use a 1-hot representation\pause
  \item Use winner-takes-all prediction (largest output)\pause
  \end{itemize}
\end{PauseHighLight}

\end{slide}

%%%%%%%%%%%%%%%%%%%%%%% Next Slide %%%%%%%%%%%%%%%%%%%%%%%

\begin{slide}
\section{Softmax}

\begin{PauseHighLight}
  \begin{itemize}
  \item A slightly better network is to use a \emph{softmax} output
    where
    \begin{align*}
      f_i(\bm{x} | \bm{\theta}) = \frac{
      \e{\sum\limits_{j=1}^h \bar{W}_{ij} h_j(\bm{x})+\bar{b}_i} }{
      \sum_{k=1}^3 \e{\sum\limits_{j=1}^h \bar{W}_{kj} h_j(\bm{x})+\bar{b}_i} }
    \end{align*}
  \item where $h_j(\bm{x})$ is the output of the $j^{th}$ hidden
    unit\pause
  \item Note that we can the interpret the output as
    \begin{align*}
      f_i(\bm{x} | \bm{\theta}) = 
      \Prob{\text{class $i$} | \bm{x},\, \bm{\theta}}\pause
    \end{align*}
  \end{itemize}
\end{PauseHighLight}

\end{slide}

%%%%%%%%%%%%%%%%%%%%%%% Next Slide %%%%%%%%%%%%%%%%%%%%%%%

\begin{slide}
\section{Relative Entropy}

\begin{PauseHighLight}
  \begin{itemize}
  \item Given a training set $\mathcal{D} = \left( (\bm{x}^\mu,
      y^\mu) \, \big\vert \, \mu  = 1,\,2,\, \ldots,\, P \right)$ where
    $y^\mu\in \mathcal{S} = \{\text{$\alpha$-helix, $\beta$-sheet,
      Coil}\}$,\pause{} the log-likelihood of the data is
      \begin{align*}
        L(\bm{\theta}) &= \logg{\Prob{\mathcal{D}|\bm{\theta}}}
        = \sum_{\mu=1}^P \logg{\Prob{y^\mu|\bm{x}^\mu,\,
                         \bm{\theta}}}\pause \\
        &= \sum_{\mu=1}^P \logg{f_{y^\mu}(\bm{x}^\mu | \bm{\theta})}\pause
      \end{align*}
    \item This is known as the \emph{cross entropy}\pause
    \item Rather than minimising the squared error we can maximise the
      relative entropy\pause---again we use a gradient based technique\pause
  \end{itemize}
\end{PauseHighLight}

\end{slide}



%%%%%%%%%%%%%%%%%%%%%%% Next Slide %%%%%%%%%%%%%%%%%%%%%%%

\begin{slide}
\section{Simple ANN}

\begin{PauseHighLight}
  \begin{itemize}
  \item A simple network would consist of
    \begin{itemize}
    \item $W\times 21$ inputs where $W$ is the window size (typically
      around 11)\pause
    \item $n_h$ hidden sigmoidal units (typically around 50)\pause
    \item 3 softmax output units\pause
    \end{itemize}
  \item Results are around 60\% prediction of the three classes\pause
  \end{itemize}
\end{PauseHighLight}

\end{slide}

%%%%%%%%%%%%%%%%%%%%%%% Next Slide %%%%%%%%%%%%%%%%%%%%%%%

\begin{slide}
\section[-2]{Learning A Better Input Representation}

\begin{PauseHighLight}
  \begin{itemize}
  \item The 1-hot vector representation of the input makes the network
    very large\pause
  \item We can learn an encoding or \emph{embedding} by using the same
    21 to $M$ layer for all inputs in the window\pause
    \begin{center}
      \includegraphics[width=\linewidth]{weightSharing}
    \end{center}
  \item All the input encoding networks share the same weights\pause
  \item This slightly complicates back-prop, but is not too difficult\pause
  \end{itemize}
\end{PauseHighLight}

\end{slide}

%%%%%%%%%%%%%%%%%%%%%%% Next Slide %%%%%%%%%%%%%%%%%%%%%%%

\begin{slide}
\section{Top Tier}

\begin{PauseHighLight}
  \begin{itemize}
  \item When testing we get a prediction at each site of the
    sequence\pause
  \item We can improve our machine by building a second machine that
    takes the prediction at a site and its neighbours\pause
  \item This uses slightly more global information and prevents silly
    predictions, e.g. \textcolor{red}{$\alpha$}, \textcolor{blue}{$\beta$},
    \textcolor{red}{$\alpha$}, \textcolor{green}{$C$}\pause
  \item In machine learning this is known as stacking\pause
  \end{itemize}
\end{PauseHighLight}

\end{slide}

%%%%%%%%%%%%%%%%%%%%%%% Next Slide %%%%%%%%%%%%%%%%%%%%%%%

\begin{slide}
\section{Using Profiles}

\begin{PauseHighLight}
  \begin{itemize}
  \item The big improvement in prediction came from using the protein
    profile obtained by finding similar protein structures
    \begin{center}
      \includegraphics[width=\linewidth]{proteinProfile}
    \end{center}
  \end{itemize}
\end{PauseHighLight}

\end{slide}

%%%%%%%%%%%%%%%%%%%%%%% Next Slide %%%%%%%%%%%%%%%%%%%%%%%

\begin{slide}
\section{Profile Encoding}

\begin{PauseHighLight}
  \begin{itemize}
  \item Instead of using the protein sequence as input we use a coding
    showing all possible substitutions, deletions or insertions at a
    site\pause
  \item This works because protein structure is much better conserved
    than protein sequences\pause
  \item Thus related protein sequences will with high probability have
    the same secondary structure\pause
  \item The protein profiles are ``smoothed'' by the related sequence
    information\pause
  \item This improved performance by $\sim5\%$\pause
  \end{itemize}
\end{PauseHighLight}

\end{slide}

%%%%%%%%%%%%%%%%%%%%%%% Next Slide %%%%%%%%%%%%%%%%%%%%%%%

\begin{slide}
\section[-1]{PSIPRED}

\begin{PauseHighLight}
  \begin{itemize}
  \item Profiling can be taken further.  Rather than just finding
    matching sequences we can iterate the process\pause
  \item We use the matched sequences to find other related
    sequences\pause
  \item This is what Psi-Blast does\pause
  \item It slow, but it produces the best profiles\pause
  \item The best predictor for many years was PSIPRED, an MLP that uses
    Psi-Blast profiles, getting close to 80\% performance\pause---now
    deep learning (CNNs, etc.) dominate\pauseb
  \end{itemize}
\end{PauseHighLight}

\end{slide}

%%%%%%%%%%%%%%%%%%%%%%% Next Slide %%%%%%%%%%%%%%%%%%%%%%%

\begin{slide}
\section{Summary}

\begin{PauseHighLight}
  \begin{itemize}
  \item Common machine learning tools have proved very useful in
    biological sequence analysis\pause
  \item However, great care must be taken to ensure that the information
    fed into them is meaningful\pause
  \item Often these have to be combined with other tools (such as
    profilers) to get really good performance\pause
  \end{itemize}
\end{PauseHighLight}

\end{slide}


